\documentclass[aspectratio=169]{beamer}
\usepackage{tikz}
\usepackage{graphicx}
\usepackage{listings}
\usepackage{xcolor}
\usepackage{hyperref}
\usepackage[spanish,es-noquoting]{babel}
\usepackage[round]{natbib}

\lstset{
  basicstyle=\ttfamily\small,
  keywordstyle=\color{blue},
  stringstyle=\color{red},
  commentstyle=\color{gray},
  showstringspaces=false
}
% Minimal look
\usetheme{Warsaw}
%\setbeamertemplate{headline}
\setbeamertemplate{footline}{
  \leavevmode%
  \hbox{%
    % Left: Section name
    \begin{beamercolorbox}[wd=.33\paperwidth,ht=2.25ex,dp=1ex,left]{author in head/foot}%
      \hspace*{2ex}\insertsection
    \end{beamercolorbox}%
    % Center: Presentation title
    \begin{beamercolorbox}[wd=.34\paperwidth,ht=2.25ex,dp=1ex,center]{title in head/foot}%
      \insertshorttitle
    \end{beamercolorbox}%
    % Right: Slide numbers
    \begin{beamercolorbox}[wd=.33\paperwidth,ht=2.25ex,dp=1ex,right]{date in head/foot}%
      \insertframenumber{} / \inserttotalframenumber\hspace*{2ex}
    \end{beamercolorbox}%
  }%
  \vskip0pt%
}

\usecolortheme{default}
\usefonttheme{serif}
\graphicspath{{../}{./}}
\setbeamertemplate{caption}[numbered]

% Title info
\title[Módulo 1: Introducción a la In]{Módulo 1: Introducción a la Inteligencia Artificial}
\subtitle{Mejorando nuestro día a día con IA}
\author{Ricardo Reyes-Jimenez}
\institute{IA para Todos}
\date{\today}



\begin{document}

% Title slide
\begin{frame}
  \titlepage
\end{frame}

% Reset background for normal slides
\usebackgroundtemplate{}



\begin{frame}[fragile]{Objetivo de la Sesión}


\textbf{Al finalizar este módulo, los participantes serán capaces de:}

1. Comprender los fundamentos de la Inteligencia Artificial en la vida cotidiana \cite{openai2023}.
2. Aplicar la IA para resolver problemas diarios en diferentes profesiones.
3. Mejorar la productividad y calidad del trabajo mediante herramientas de IA \cite{mckinsey2023}.

\end{frame}



\begin{frame}[fragile]{1. La IA en el Consultorio (Médicos)}


\textbf{La IA no reemplaza el criterio médico, pero acelera el diagnóstico y la investigación.}

\begin{itemize}
  \item \textbf{Caso Práctico:} Un médico necesita resumir el historial clínico de un paciente complejo o buscar interacciones farmacológicas recientes.
  \item \textbf{Uso de IA:} ``Actúa como un asistente médico experto. Resume los siguientes 5 años de historial clínico, destacando alergias y condiciones crónicas, y enumera las posibles interacciones entre el Medicamento A y el Medicamento B según la literatura reciente.''
\end{itemize}
\textbf{Resultado:} Más tiempo para la atención al paciente y menos tiempo en el papeleo.

\end{frame}



\begin{frame}[fragile]{2. Optimizando Cálculos y Estructuras (Ingenieros)}


\textbf{Para los ingenieros, la IA es una herramienta de validación y automatización.}

\begin{itemize}
  \item \textbf{Caso Práctico:} Un ingeniero civil necesita redactar un informe sobre la fatiga de materiales para un cliente no técnico.
  \item \textbf{Uso de IA:} ``Actúa como un ingeniero civil senior. Explica los riesgos de la fatiga de materiales en puentes colgantes a una audiencia sin conocimientos técnicos. Usa metáforas sencillas y mantén el tono profesional e informativo.''
\end{itemize}
\textbf{Resultado:} Comunicación clara con los clientes y reportes más rápidos.

\end{frame}



\begin{frame}[fragile]{3. Análisis de Mercados y Tendencias (Economistas)}


\textbf{La IA procesa grandes volúmenes de datos económicos en segundos.}

\begin{itemize}
  \item \textbf{Caso Práctico:} Un economista quiere analizar el sentimiento del mercado basado en las últimas noticias financieras.
  \item \textbf{Uso de IA:} ``Actúa como un analista financiero. Toma el siguiente texto de noticias económicas y extrae el sentimiento general del mercado. Identifica las 3 principales tendencias y preséntalas en formato de viñetas (bullet points).''
\end{itemize}
\textbf{Resultado:} Decisiones más rápidas basadas en resúmenes precisos de la información.

\end{frame}



\begin{frame}[fragile]{4. Eficiencia en el Trabajo Diario (Empleados)}


\textbf{Cualquier oficinista puede reducir horas de trabajo repetitivo \cite{brynjolfsson2023}.}

\begin{itemize}
  \item \textbf{Caso Práctico:} Un empleado necesita redactar un correo electrónico delicado a su jefe solicitando una extensión de tiempo para un proyecto.
  \item \textbf{Uso de IA:} ``Necesito escribir un correo a mi director. \textbf{Contexto:} El proyecto se retrasó 2 días por problemas con el proveedor. \textbf{Audiencia:} Mi jefe, que prefiere soluciones directas. \textbf{Método:} Un correo breve y profesional que empiece con el problema, seguido inmediatamente por mi plan para resolverlo trabajando unas horas extras el fin de semana.''
\end{itemize}
\textbf{Resultado:} Correos efectivos, profesionales y con menor nivel de estrés al redactarlos.

\end{frame}



\begin{frame}[fragile]{Recapitulando el Módulo}


\textbf{Hoy hemos aprendido:}
\begin{itemize}
  \item La IA es una herramienta versátil para cualquier profesión.
  \item Contextualizar el problema (rol, tarea, audiencia) mejora drásticamente los resultados.
  \item La automatización de tareas tediosas nos permite enfocarnos en el trabajo de alto valor.
\end{itemize}
\textbf{¡Sigan aprendiendo, sigan trabajando y sigan disfrutando de la Inteligencia Artificial!}
\end{frame}




\begin{frame}[allowframebreaks]{References}
  \bibliographystyle{apalike}
  \bibliography{references}
\end{frame}


\end{document}